% GENERATED FILE. Edit canonical lessons, then run npm run generate:books.
% canonical-source-hash: fnv1a64:6cce3decfbdaad8b
% canonical-lessons: TE-C528-tobuttuvu, TE-C528-annadammulu, TE-C528-maridi, TE-C528-adapaducu, TE-C528-vadina

\chapter{Brothers and Sisters}
\label{ch:te-brothers-and-sisters}
\hlchaptermodality{528}

\begin{chapteropening}
\textbf{By the end of this chapter:} \emph{I can talk about my siblings and my brothers' and husband's family.}

Complete the last lesson of chapter 528: I can talk about my siblings and my brothers' and husband's family.
\end{chapteropening}

\section[\texorpdfstring{\te{తోబుట్టువు} (tōbuṭṭuvu)}{తోబుట్టువు (tōbuṭṭuvu)}]{\te{తోబుట్టువు} (tōbuṭṭuvu) — a sibling}

\label{lesson:TE-C528-tobuttuvu}



\begin{quote}
Before the new one: say the Telugu for a resignation, then the Telugu for to quit.
\end{quote}

\subsection*{You'll want to know: \te{తోబుట్టువు}}
\textbf{\te{తోబుట్టువు}} — \emph{tōbuṭṭuvu} — \textquotedblleft{}a sibling\textquotedblright{}.

\begin{quote}
\te{నాకు} \te{ఇద్దరు} \te{తోబుట్టువులు} \te{ఉన్నారు}.
\end{quote}

\begin{quote}
\emph{nāku iddaru tōbuṭṭuvulu unnāru.}
\end{quote}

\begin{quote}
I have two siblings.
\end{quote}

Plural: \te{తోబుట్టువులు}, \emph{tōbuṭṭuvulu}.

Literally born with: brothers and sisters alike.

\begin{cousinweb}
\textbf{\te{తోబుట్టువు}} (\emph{tōbuṭṭuvu}) is a native compound: \emph{tōḍa}, \textquotedblleft{}with\textquotedblright{} (DEDR 3563), and \emph{puṭṭu}, \textquotedblleft{}to be born\textquotedblright{} (DEDR 4264): one born with you.
\end{cousinweb}

\begin{grammarlens}[title={What you've built}]
One more word for describing what you lived and felt.
\end{grammarlens}

\subsection*{Your turn}
\begin{itemize}
\raggedright
  \item \emph{Say it:} \emph{tōbuṭṭuvu}
  \item \emph{Say it:} \emph{tōbuṭṭuvu}, once more
  \item \emph{Say it:} \emph{nāku iddaru tōbuṭṭuvulu unnāru.} — I have two siblings.
  \item \emph{Recall:} say the Telugu for a resignation, then the Telugu for to quit, then say \emph{tōbuṭṭuvu} again
\end{itemize}

\begin{tcolorbox}[breakable,colback=teal!4,colframe=teal!35!black,title={Before you move on}]
What does \te{తోబుట్టువు} mean? (\textquotedblleft{}A sibling\textquotedblright{}.) Say it once more.
\end{tcolorbox}
\section[\texorpdfstring{\te{అన్నదమ్ములు} (annadammulu)}{అన్నదమ్ములు (annadammulu)}]{\te{అన్నదమ్ములు} (annadammulu) — brothers; a set of brothers}

\label{lesson:TE-C528-annadammulu}



\begin{quote}
Before the new one: say the Telugu for to quit, then the Telugu for a sibling.
\end{quote}

\subsection*{You'll want to know: \te{అన్నదమ్ములు}}
\textbf{\te{అన్నదమ్ములు}} — \emph{annadammulu} — \textquotedblleft{}brothers; a set of brothers\textquotedblright{}.

\begin{quote}
\te{ఆ} \te{అన్నదమ్ములు} \te{కలిసి} \te{వ్యాపారం} \te{చేస్తున్నారు}.
\end{quote}

\begin{quote}
\emph{ā annadammulu kalisi vyāpāraṁ cēstunnāru.}
\end{quote}

\begin{quote}
Those brothers run a business together.
\end{quote}

A pair word: \te{అన్న}, elder brother, + \te{తమ్ముడు}, younger brother.

\begin{cousinweb}
\textbf{\te{అన్నదమ్ములు}} (\emph{annadammulu}) is a native pair word: \emph{anna}, \textquotedblleft{}elder brother\textquotedblright{} (DEDR 131), and \emph{tammuḍu}, \textquotedblleft{}younger brother\textquotedblright{} (DEDR 3085).
\end{cousinweb}

\begin{grammarlens}[title={What you've built}]
One more word for describing what you lived and felt.
\end{grammarlens}

\subsection*{Your turn}
\begin{itemize}
\raggedright
  \item \emph{Say it:} \emph{annadammulu}
  \item \emph{Say it:} \emph{annadammulu}, once more
  \item \emph{Say it:} \emph{ā annadammulu kalisi vyāpāraṁ cēstunnāru.} — Those brothers run a business together.
  \item \emph{Recall:} say the Telugu for to quit, then the Telugu for a sibling, then say \emph{annadammulu} again
\end{itemize}

\begin{tcolorbox}[breakable,colback=teal!4,colframe=teal!35!black,title={Before you move on}]
What does \te{అన్నదమ్ములు} mean? (\textquotedblleft{}Brothers; a set of brothers\textquotedblright{}.) Say it once more.
\end{tcolorbox}
\section[\texorpdfstring{\te{మరిది} (maridi)}{మరిది (maridi)}]{\te{మరిది} (maridi) — a husband's younger brother}

\label{lesson:TE-C528-maridi}



\begin{quote}
Before the new one: say the Telugu for a sibling, then the Telugu for brothers.
\end{quote}

\subsection*{You'll want to know: \te{మరిది}}
\textbf{\te{మరిది}} — \emph{maridi} — \textquotedblleft{}a husband's younger brother\textquotedblright{}.

\begin{quote}
\te{పెళ్ళి} \te{తరువాత} \te{మా} \te{మరిది} \te{మాతోనే} \te{ఉన్నాడు}.
\end{quote}

\begin{quote}
\emph{peḷḷi taruvāta mā maridi mātōnē unnāḍu.}
\end{quote}

\begin{quote}
After the wedding my husband's younger brother lived with us.
\end{quote}

Also a younger sister's husband.

\begin{cousinweb}
\textbf{\te{మరిది}} (\emph{maridi}) is native Dravidian: DEDR 4762 lists Telugu \emph{maṟãdi}, \textquotedblleft{}a spouse's younger brother, a younger sister's husband\textquotedblright{}, beside \emph{maṟãdalu}, the same for women.
\end{cousinweb}

\begin{grammarlens}[title={What you've built}]
One more word for describing what you lived and felt.
\end{grammarlens}

\subsection*{Your turn}
\begin{itemize}
\raggedright
  \item \emph{Say it:} \emph{maridi}
  \item \emph{Say it:} \emph{maridi}, once more
  \item \emph{Say it:} \emph{peḷḷi taruvāta mā maridi mātōnē unnāḍu.} — After the wedding my husband's younger brother lived with us.
  \item \emph{Recall:} say the Telugu for a sibling, then the Telugu for brothers, then say \emph{maridi} again
\end{itemize}

\begin{tcolorbox}[breakable,colback=teal!4,colframe=teal!35!black,title={Before you move on}]
What does \te{మరిది} mean? (\textquotedblleft{}A husband's younger brother\textquotedblright{}.) Say it once more.
\end{tcolorbox}
\section[\texorpdfstring{\te{ఆడపడుచు} (āḍapaḍucu)}{ఆడపడుచు (āḍapaḍucu)}]{\te{ఆడపడుచు} (āḍapaḍucu) — a husband's sister}

\label{lesson:TE-C528-adapaducu}



\begin{quote}
Before the new one: say the Telugu for brothers, then the Telugu for a husband's younger brother.
\end{quote}

\subsection*{You'll want to know: \te{ఆడపడుచు}}
\textbf{\te{ఆడపడుచు}} — \emph{āḍapaḍucu} — \textquotedblleft{}a husband's sister\textquotedblright{}.

\begin{quote}
\te{పండుగకు} \te{ఆడపడుచు} \te{ఇంటికి} \te{వచ్చింది}.
\end{quote}

\begin{quote}
\emph{paṇḍugaku āḍapaḍucu iṇṭiki vaccindi.}
\end{quote}

\begin{quote}
My husband's sister came home for the festival.
\end{quote}

Also a married daughter of the house visiting home.

\begin{cousinweb}
\textbf{\te{ఆడపడుచు}} (\emph{āḍapaḍucu}) is built from \te{ఆడ}, \textquotedblleft{}female\textquotedblright{}, and native \emph{paḍucu}, \textquotedblleft{}a girl\textquotedblright{}, which DEDR 3840 lists with \emph{paḍãti}, \textquotedblleft{}a woman\textquotedblright{}.
\end{cousinweb}

\begin{grammarlens}[title={What you've built}]
One more word for describing what you lived and felt.
\end{grammarlens}

\subsection*{Your turn}
\begin{itemize}
\raggedright
  \item \emph{Say it:} \emph{āḍapaḍucu}
  \item \emph{Say it:} \emph{āḍapaḍucu}, once more
  \item \emph{Say it:} \emph{paṇḍugaku āḍapaḍucu iṇṭiki vaccindi.} — My husband's sister came home for the festival.
  \item \emph{Recall:} say the Telugu for brothers, then the Telugu for a husband's younger brother, then say \emph{āḍapaḍucu} again
\end{itemize}

\begin{tcolorbox}[breakable,colback=teal!4,colframe=teal!35!black,title={Before you move on}]
What does \te{ఆడపడుచు} mean? (\textquotedblleft{}A husband's sister\textquotedblright{}.) Say it once more.
\end{tcolorbox}
\section[\texorpdfstring{\te{వదిన} (vadina)}{వదిన (vadina)}]{\te{వదిన} (vadina) — an elder brother's wife}

\label{lesson:TE-C528-vadina}



\begin{quote}
Before the new one: say the Telugu for a husband's younger brother, then the Telugu for a husband's sister.
\end{quote}

\subsection*{You'll want to know: \te{వదిన}}
\textbf{\te{వదిన}} — \emph{vadina} — \textquotedblleft{}an elder brother's wife\textquotedblright{}.

\begin{quote}
\te{మా} \te{వదిన} \te{చాలా} \te{బాగా} \te{వంట} \te{చేస్తుంది}.
\end{quote}

\begin{quote}
\emph{mā vadina cālā bāgā vaṇṭa cēstundi.}
\end{quote}

\begin{quote}
My elder brother's wife cooks very well.
\end{quote}

Also a husband's or a wife's elder sister.

\begin{cousinweb}
\textbf{\te{వదిన}} (\emph{vadina}) is native Dravidian: DEDR 4690 lists Telugu \emph{vadine}, \textquotedblleft{}an elder sister-in-law, an elder brother's wife\textquotedblright{}.
\end{cousinweb}

\begin{grammarlens}[title={What you've built}]
One more word for describing what you lived and felt.
\end{grammarlens}

\subsection*{Your turn}
\begin{itemize}
\raggedright
  \item \emph{Say it:} \emph{vadina}
  \item \emph{Say it:} \emph{vadina}, once more
  \item \emph{Say it:} \emph{mā vadina cālā bāgā vaṇṭa cēstundi.} — My elder brother's wife cooks very well.
  \item \emph{Recall:} say the Telugu for a husband's younger brother, then the Telugu for a husband's sister, then say \emph{vadina} again
\end{itemize}

\begin{tcolorbox}[breakable,colback=teal!4,colframe=teal!35!black,title={Before you move on}]
What does \te{వదిన} mean? (\textquotedblleft{}An elder brother's wife\textquotedblright{}.) Say it once more.
\end{tcolorbox}
